\chapter{尚未解决的问题与两类方法的能力边界}\label{ch:open}
\section{哪些问题在方程层面就丢失了信息}
有限矩模型只知道分布函数的若干积分。不同的 $f_s$ 可以具有相同密度、平均速度和压力张量，却有完全不同的高能尾、束流和俘获粒子。因此五矩／十矩系统本身不能唯一恢复能谱、非热加速效率、速度空间精细结构或每个粒子的轨道历史。用测试粒子补充时，若高能粒子携带显著压力或电流，还必须反作用到流体，否则能量和动力学反馈不闭合。

普通Hall MHD不保留离子完整有限回旋半径、离子多流、电子非旋转对称压力或无碰撞热流的完整信息；因此当问题恰好由这些机制决定时，单纯提高网格分辨率不能修复缺失的方程。十矩在压力层面前进一步，但三阶和更高矩仍由闭合猜测。

非相对论电子离子双流体还不能直接回答辐射反作用、对产生、强相对论焓、部分电离化学或多重电荷态问题。检索中发现这些邻近文献时应保留索引，但不把它们与本文主体模型混为同一方程组。

\section{哪些问题可能用Hall／双流体有效研究}
\subsection{从单X线到湍流态的相界}
基于 \citet{Huang2011,Huang2024}，可扫描 $S,L/d_i,L_z/d_i,B_g/B_0$，测量长时间单X线占空比、磁通绳数、净通量率与三维谱。尚不明确的是状态是否由少数无量纲参数统一控制，还是电子闭合和边界也不可约。可先用Hall模型建立相图，再挑相界附近少量点用十矩和PIC比较。

判别失败的条件要预先设定：若改变超电阻或电子质量后相界迁移幅度与扫描范围同阶，则应称为该闭合的相图，不能称为自然等离子体的普适相界。

\subsection{电子层长度到底由谁选定}
竞争机制包括外部驱动、下游背压、电子惯性、压力张量、黏性／超电阻、电子剪切不稳定性及三维破裂。建议用同一上游和边界，分别启用电阻、惯性五矩、十矩局地闭合、十矩非局地闭合。核心输出不是单一速率，而是 $L_e/d_i$、$\delta_e/d_e$、电子动量通量和欧姆定律分解。\citep{Sullivan2009,Ng2015,Ng2017}

\subsection{触发时间能否预测而不是被网格决定}
GEM强种子比较不能独自解答自然触发。需从宏观宽片开始、控制驱动和噪声，并测量到达非理想尺度的时间。\citet{Arshad2026}的全局MHD--PIC比较发现其设置下理想及Hall MHD触发具有网格依赖；这对Hall模型的警示是必须明确电子非理想闭合。该结果不等于所有有物理电子闭合的Hall模型都永远无法收敛，也不能把该文报告的PIC网格要求直接移植为流体通用准则。

\subsection{压力闭合是否可跨问题迁移}
在Harris片校准的 $k_0$ 能否用于岛合并、不对称片、强导引场和三维低混合不稳定性？这是十矩的关键开放问题。可以采用“训练／验证分开”的设计：一个算例定参数，另外三个算例不再调参；比较压力张量、热流、增长率和混合量。对数据驱动闭合还需强制旋转协变、能量一致性及压力正定性，避免只拟合重联率。\citep{Wang2015,Ng2020,Closure2026}

\subsection{峰值率、平均率和全局时间的统一}
在岛合并中，局地快率与整体合并慢并存。可以用系统尺寸扫描建立 $\max\R$、$\langle\R\rangle$ 与合并时间的联合标度，检查上游通量供应和反复重联是否解释差异。Hall模型计算成本相对低，适合做这一问题的参数初筛，但离子动力学敏感区域要借助十矩／混合模型比较。\citep{Knoll2006,Ng2015}

\section{值得直接执行的三个研究方案}
\begin{longtable}{P{22mm}P{44mm}P{48mm}P{35mm}}
\toprule 方案 & 参数与控制 & 需要观察的新结果 & 可信结论条件\\\midrule
A：闭合改变三维形态 & 固定 $m_i/m_e,c/\VA$；局地／梯度／非局地闭合对照；扫描 $L_z$ & 横向不稳定性阈值、混合、磁通绳寿命 & 各闭合独立网格收敛；热流与压力证据\\
B：磁岛触发与迟滞 & 同时扫描 $S,L/d_i$；快态慢态两种初值 & 单X线与多X点切换、迟滞随三维变化 & 扫描速率、等待时间与边界独立检验\\
C：实验几何的下游约束 & 局地平板对照MRX／合并几何；固定上游驱动 & 四极场畸变、片漂移、全局时间变化 & 使用实际边界、合成诊断与参数误差条\\\bottomrule
\end{longtable}
这三个方案是本报告基于文献差异提出的研究建议，不是已经得到的新科学发现。建议A最能检验双流体闭合价值，建议B最适合发现不同渐近状态，建议C最方便与实验观测联系。

\section{“没有解决”也要分层表述}
\begin{enumerate}
\item \textbf{原则上信息不足：}仅从有限矩恢复任意非热分布，必须增加模型。
\item \textbf{模型能够表达但闭合未可靠：}非对角压力、无碰撞热流、部分次级不稳定性。
\item \textbf{模型能算但参数范围不足：}真实质量比、大 $L/d_i$、三维长时间统计。
\item \textbf{观测尚不能唯一判别：}某一四极场、单次喷流或热化峰值的机制归属。
\end{enumerate}
这四类问题对应增加物理、改进闭合、提升计算和改进测量四种不同工作。不能把它们全部称为“还需要更高分辨率”。
